\documentclass[10pt,twocolumn]{article}
\usepackage[margin=0.8in]{geometry}
\usepackage{graphicx}
\usepackage{booktabs}
\usepackage{amsmath}
\usepackage{hyperref}
\usepackage{natbib}
% generated by analysis/make_claims.py
\expandafter\def\csname res@band_hi\endcsname{0.8}
\expandafter\def\csname res@band_lo\endcsname{0.2}
\expandafter\def\csname res@c2_conf_h1n\endcsname{24}
\expandafter\def\csname res@c2_conf_h2p\endcsname{0.043}
\expandafter\def\csname res@c2_haiku_all_gate\endcsname{passed}
\expandafter\def\csname res@c2_haiku_all_h1ci90\endcsname{[+0.04, +0.19]}
\expandafter\def\csname res@c2_haiku_all_h1d\endcsname{+0.12}
\expandafter\def\csname res@c2_haiku_all_h1eq\endcsname{not inside}
\expandafter\def\csname res@c2_haiku_all_h1n\endcsname{103}
\expandafter\def\csname res@c2_haiku_all_h2\endcsname{+0.12}
\expandafter\def\csname res@c2_haiku_all_h2ci\endcsname{[+0.04, +0.19]}
\expandafter\def\csname res@c2_haiku_all_h2p\endcsname{0.001}
\expandafter\def\csname res@c2_haiku_all_isowrong\endcsname{0.15}
\expandafter\def\csname res@c2_haiku_all_p1\endcsname{0.007}
\expandafter\def\csname res@c2_haiku_all_solve\endcsname{0.28}
\expandafter\def\csname res@c2_haiku_all_ta_iso\endcsname{0.30}
\expandafter\def\csname res@c2_haiku_all_ta_pick\endcsname{0.42}
\expandafter\def\csname res@c2_haiku_conf_gate\endcsname{failed}
\expandafter\def\csname res@c2_haiku_conf_h1ci90\endcsname{[-0.12, +0.12]}
\expandafter\def\csname res@c2_haiku_conf_h1d\endcsname{+0.00}
\expandafter\def\csname res@c2_haiku_conf_h1eq\endcsname{not inside}
\expandafter\def\csname res@c2_haiku_conf_h1n\endcsname{24}
\expandafter\def\csname res@c2_haiku_conf_h2\endcsname{+0.12}
\expandafter\def\csname res@c2_haiku_conf_h2ci\endcsname{[+0.00, +0.29]}
\expandafter\def\csname res@c2_haiku_conf_h2p\endcsname{0.043}
\expandafter\def\csname res@c2_haiku_conf_isowrong\endcsname{0.04}
\expandafter\def\csname res@c2_haiku_conf_solve\endcsname{0.88}
\expandafter\def\csname res@c2_haiku_conf_ta_iso\endcsname{0.12}
\expandafter\def\csname res@c2_haiku_conf_ta_pick\endcsname{0.12}
\expandafter\def\csname res@c2_haiku_repl_gate\endcsname{failed}
\expandafter\def\csname res@c2_haiku_repl_h1ci90\endcsname{[+0.06, +0.24]}
\expandafter\def\csname res@c2_haiku_repl_h1d\endcsname{+0.15}
\expandafter\def\csname res@c2_haiku_repl_h1eq\endcsname{not inside}
\expandafter\def\csname res@c2_haiku_repl_h1n\endcsname{79}
\expandafter\def\csname res@c2_haiku_repl_h2\endcsname{+0.11}
\expandafter\def\csname res@c2_haiku_repl_h2ci\endcsname{[+0.03, +0.20]}
\expandafter\def\csname res@c2_haiku_repl_h2p\endcsname{0.006}
\expandafter\def\csname res@c2_haiku_repl_isowrong\endcsname{0.18}
\expandafter\def\csname res@c2_haiku_repl_solve\endcsname{0.10}
\expandafter\def\csname res@c2_haiku_repl_ta_iso\endcsname{0.35}
\expandafter\def\csname res@c2_haiku_repl_ta_pick\endcsname{0.51}
\expandafter\def\csname res@c2_holm_haiku_H1\endcsname{0.621}
\expandafter\def\csname res@c2_holm_haiku_H2\endcsname{0.005}
\expandafter\def\csname res@c2_holm_sonnet_H1\endcsname{0.366}
\expandafter\def\csname res@c2_holm_sonnet_H2\endcsname{0.366}
\expandafter\def\csname res@c2_repl_h2\endcsname{+0.11}
\expandafter\def\csname res@c2_sonnet_conf_gate\endcsname{failed}
\expandafter\def\csname res@c2_sonnet_conf_h1ci90\endcsname{[-0.12, +0.12]}
\expandafter\def\csname res@c2_sonnet_conf_h1d\endcsname{+0.00}
\expandafter\def\csname res@c2_sonnet_conf_h1eq\endcsname{not inside}
\expandafter\def\csname res@c2_sonnet_conf_h1n\endcsname{17}
\expandafter\def\csname res@c2_sonnet_conf_h2\endcsname{+0.11}
\expandafter\def\csname res@c2_sonnet_conf_h2ci\endcsname{[+0.00, +0.26]}
\expandafter\def\csname res@c2_sonnet_conf_h2p\endcsname{0.122}
\expandafter\def\csname res@c2_sonnet_conf_isowrong\endcsname{0.03}
\expandafter\def\csname res@c2_sonnet_conf_solve\endcsname{0.92}
\expandafter\def\csname res@c2_sonnet_conf_ta_iso\endcsname{0.12}
\expandafter\def\csname res@c2_sonnet_conf_ta_pick\endcsname{0.12}
\expandafter\def\csname res@c2_sonnet_isoparse\endcsname{0.896}
\expandafter\def\csname res@c3_disagree\endcsname{22}
\expandafter\def\csname res@c3_gate_diff\endcsname{0.10}
\expandafter\def\csname res@c3_gate_parse\endcsname{0.90}
\expandafter\def\csname res@c3_n_mixed\endcsname{50}
\expandafter\def\csname res@c3_n_solved\endcsname{137}
\expandafter\def\csname res@c3_n_ta\endcsname{278}
\expandafter\def\csname res@c3_n_unresolved\endcsname{1}
\expandafter\def\csname res@c3_n_unsolved\endcsname{90}
\expandafter\def\csname res@c3h_b\endcsname{+0.12}
\expandafter\def\csname res@c3h_b_ci\endcsname{[-0.03, +0.27]}
\expandafter\def\csname res@c3h_b_p\endcsname{0.065}
\expandafter\def\csname res@c3h_cost\endcsname{10.7}
\expandafter\def\csname res@c3h_ds\endcsname{+0.06}
\expandafter\def\csname res@c3h_ds_ci\endcsname{[+0.01, +0.11]}
\expandafter\def\csname res@c3h_ds_n\endcsname{137}
\expandafter\def\csname res@c3h_ds_p\endcsname{0.019}
\expandafter\def\csname res@c3h_du\endcsname{+0.24}
\expandafter\def\csname res@c3h_du_ci\endcsname{[+0.13, +0.36]}
\expandafter\def\csname res@c3h_du_n\endcsname{90}
\expandafter\def\csname res@c3h_du_p\endcsname{$<$0.001}
\expandafter\def\csname res@c3h_fam_arxiv\endcsname{+0.12}
\expandafter\def\csname res@c3h_fam_arxiv_ci\endcsname{[-0.04, +0.28]}
\expandafter\def\csname res@c3h_fam_arxiv_n\endcsname{140}
\expandafter\def\csname res@c3h_fam_arxiv_nu\endcsname{84}
\expandafter\def\csname res@c3h_fam_contest\endcsname{+0.11}
\expandafter\def\csname res@c3h_fam_contest_ci\endcsname{[-0.44, +0.66]}
\expandafter\def\csname res@c3h_fam_contest_n\endcsname{87}
\expandafter\def\csname res@c3h_fam_contest_nu\endcsname{6}
\expandafter\def\csname res@c3h_gate\endcsname{passed}
\expandafter\def\csname res@c3h_h2s\endcsname{+0.02}
\expandafter\def\csname res@c3h_h2s_ci\endcsname{[-0.01, +0.05]}
\expandafter\def\csname res@c3h_h2s_n\endcsname{137}
\expandafter\def\csname res@c3h_h2s_p\endcsname{0.128}
\expandafter\def\csname res@c3h_h2u\endcsname{+0.19}
\expandafter\def\csname res@c3h_h2u_ci\endcsname{[+0.11, +0.28]}
\expandafter\def\csname res@c3h_h2u_n\endcsname{90}
\expandafter\def\csname res@c3h_h2u_p\endcsname{$<$0.001}
\expandafter\def\csname res@c3h_holm_H1\endcsname{0.129}
\expandafter\def\csname res@c3h_holm_H1u\endcsname{$<$0.001}
\expandafter\def\csname res@c3h_holm_H2s\endcsname{0.129}
\expandafter\def\csname res@c3h_holm_H2u\endcsname{$<$0.001}
\expandafter\def\csname res@c3h_isoparse\endcsname{1.000}
\expandafter\def\csname res@c3h_isoright\endcsname{0.83}
\expandafter\def\csname res@c3h_isowrong\endcsname{0.08}
\expandafter\def\csname res@c3h_n\endcsname{227}
\expandafter\def\csname res@c3h_solved_iso\endcsname{0.06}
\expandafter\def\csname res@c3h_solved_pick\endcsname{0.12}
\expandafter\def\csname res@c3h_unadj\endcsname{+0.19}
\expandafter\def\csname res@c3h_unadj_ci\endcsname{[+0.06, +0.31]}
\expandafter\def\csname res@c3h_unsolved_iso\endcsname{0.36}
\expandafter\def\csname res@c3h_unsolved_pick\endcsname{0.60}
\expandafter\def\csname res@c3h_verdict_H1\endcsname{inconclusive}
\expandafter\def\csname res@c3h_verdict_H1u\endcsname{confirmed}
\expandafter\def\csname res@c3h_verdict_H2s\endcsname{inconclusive}
\expandafter\def\csname res@c3h_verdict_H2u\endcsname{confirmed}
\expandafter\def\csname res@c3s_b\endcsname{-0.02}
\expandafter\def\csname res@c3s_b_ci\endcsname{[-0.15, +0.11]}
\expandafter\def\csname res@c3s_b_p\endcsname{0.614}
\expandafter\def\csname res@c3s_cost\endcsname{66.2}
\expandafter\def\csname res@c3s_ds\endcsname{+0.06}
\expandafter\def\csname res@c3s_ds_ci\endcsname{[+0.01, +0.11]}
\expandafter\def\csname res@c3s_ds_n\endcsname{135}
\expandafter\def\csname res@c3s_ds_p\endcsname{0.021}
\expandafter\def\csname res@c3s_du\endcsname{+0.13}
\expandafter\def\csname res@c3s_du_ci\endcsname{[+0.04, +0.23]}
\expandafter\def\csname res@c3s_du_n\endcsname{90}
\expandafter\def\csname res@c3s_du_p\endcsname{0.004}
\expandafter\def\csname res@c3s_fam_arxiv\endcsname{-0.03}
\expandafter\def\csname res@c3s_fam_arxiv_ci\endcsname{[-0.17, +0.12]}
\expandafter\def\csname res@c3s_fam_arxiv_n\endcsname{140}
\expandafter\def\csname res@c3s_fam_arxiv_nu\endcsname{84}
\expandafter\def\csname res@c3s_fam_contest\endcsname{+0.05}
\expandafter\def\csname res@c3s_fam_contest_ci\endcsname{[-0.09, +0.18]}
\expandafter\def\csname res@c3s_fam_contest_n\endcsname{85}
\expandafter\def\csname res@c3s_fam_contest_nu\endcsname{6}
\expandafter\def\csname res@c3s_gate\endcsname{passed}
\expandafter\def\csname res@c3s_h2s\endcsname{+0.03}
\expandafter\def\csname res@c3s_h2s_ci\endcsname{[-0.01, +0.06]}
\expandafter\def\csname res@c3s_h2s_n\endcsname{135}
\expandafter\def\csname res@c3s_h2s_p\endcsname{0.076}
\expandafter\def\csname res@c3s_h2u\endcsname{+0.28}
\expandafter\def\csname res@c3s_h2u_ci\endcsname{[+0.18, +0.38]}
\expandafter\def\csname res@c3s_h2u_n\endcsname{90}
\expandafter\def\csname res@c3s_h2u_p\endcsname{$<$0.001}
\expandafter\def\csname res@c3s_holm_H1\endcsname{0.614}
\expandafter\def\csname res@c3s_holm_H1u\endcsname{0.012}
\expandafter\def\csname res@c3s_holm_H2s\endcsname{0.152}
\expandafter\def\csname res@c3s_holm_H2u\endcsname{$<$0.001}
\expandafter\def\csname res@c3s_isoparse\endcsname{1.000}
\expandafter\def\csname res@c3s_isoright\endcsname{0.86}
\expandafter\def\csname res@c3s_isowrong\endcsname{0.11}
\expandafter\def\csname res@c3s_n\endcsname{225}
\expandafter\def\csname res@c3s_solved_iso\endcsname{0.08}
\expandafter\def\csname res@c3s_solved_pick\endcsname{0.14}
\expandafter\def\csname res@c3s_unadj\endcsname{+0.07}
\expandafter\def\csname res@c3s_unadj_ci\endcsname{[-0.03, +0.18]}
\expandafter\def\csname res@c3s_unsolved_iso\endcsname{0.50}
\expandafter\def\csname res@c3s_unsolved_pick\endcsname{0.63}
\expandafter\def\csname res@c3s_verdict_H1\endcsname{inconclusive}
\expandafter\def\csname res@c3s_verdict_H1u\endcsname{confirmed}
\expandafter\def\csname res@c3s_verdict_H2s\endcsname{inconclusive}
\expandafter\def\csname res@c3s_verdict_H2u\endcsname{confirmed}
\expandafter\def\csname res@cal_haiku\endcsname{0.98}
\expandafter\def\csname res@cal_haiku_nowork\endcsname{0.97}
\expandafter\def\csname res@cal_sonnet\endcsname{0.93}
\expandafter\def\csname res@cal_sonnet_nowork\endcsname{0.95}
\expandafter\def\csname res@cli_version\endcsname{2.1.317}
\expandafter\def\csname res@cost_subjects\endcsname{220}
\expandafter\def\csname res@eq_margin\endcsname{0.10}
\expandafter\def\csname res@h3_haiku_acc_haiku\endcsname{0.14}
\expandafter\def\csname res@h3_haiku_acc_sonnet\endcsname{0.16}
\expandafter\def\csname res@h3_haiku_ci\endcsname{[-0.04, +0.07]}
\expandafter\def\csname res@h3_haiku_d\endcsname{+0.02}
\expandafter\def\csname res@h3_haiku_n\endcsname{225}
\expandafter\def\csname res@h3_haiku_p\endcsname{0.294}
\expandafter\def\csname res@h3_sonnet_acc_haiku\endcsname{0.13}
\expandafter\def\csname res@h3_sonnet_acc_sonnet\endcsname{0.19}
\expandafter\def\csname res@h3_sonnet_ci\endcsname{[-0.01, +0.10]}
\expandafter\def\csname res@h3_sonnet_d\endcsname{+0.04}
\expandafter\def\csname res@h3_sonnet_n\endcsname{225}
\expandafter\def\csname res@h3_sonnet_p\endcsname{0.049}
\expandafter\def\csname res@h3_verdict\endcsname{inconclusive}
\expandafter\def\csname res@ledger_total\endcsname{74}
\expandafter\def\csname res@lvl90\endcsname{90\%}
\expandafter\def\csname res@lvl95\endcsname{95\%}
\expandafter\def\csname res@p1_gemma2_pickiso\endcsname{0.38}
\expandafter\def\csname res@p1_gemma2_pickiso_ci\endcsname{[+0.17, +0.59]}
\expandafter\def\csname res@p1_haiku_isowrong\endcsname{0.00}
\expandafter\def\csname res@p1_haiku_pick\endcsname{0.00}
\expandafter\def\csname res@p1_haiku_solve\endcsname{1.00}
\expandafter\def\csname res@p1_opus_isowrong\endcsname{0.00}
\expandafter\def\csname res@p1_opus_pick\endcsname{0.00}
\expandafter\def\csname res@p1_opus_solve\endcsname{1.00}
\expandafter\def\csname res@p1_sonnet_isowrong\endcsname{0.00}
\expandafter\def\csname res@p1_sonnet_pick\endcsname{0.00}
\expandafter\def\csname res@p1_sonnet_solve\endcsname{1.00}
\expandafter\def\csname res@p2_haiku_auroc\endcsname{0.64}
\expandafter\def\csname res@p2_sonnet_auroc\endcsname{0.59}
\expandafter\def\csname res@ph_all_ci\endcsname{[-0.01, +0.13]}
\expandafter\def\csname res@ph_all_d\endcsname{+0.06}
\expandafter\def\csname res@ph_solved_ci\endcsname{[-0.14, +0.05]}
\expandafter\def\csname res@ph_solved_d\endcsname{-0.04}
\expandafter\def\csname res@ph_solved_iso\endcsname{0.18}
\expandafter\def\csname res@ph_solved_n\endcsname{73}
\expandafter\def\csname res@ph_solved_pick\endcsname{0.14}
\expandafter\def\csname res@ph_unsolved_ci\endcsname{[+0.04, +0.25]}
\expandafter\def\csname res@ph_unsolved_d\endcsname{+0.15}
\expandafter\def\csname res@ph_unsolved_iso\endcsname{0.39}
\expandafter\def\csname res@ph_unsolved_n\endcsname{89}
\expandafter\def\csname res@ph_unsolved_pick\endcsname{0.54}
\expandafter\def\csname res@s10_hetero\endcsname{0.54}
\expandafter\def\csname res@s10_homo\endcsname{0.67}
\expandafter\def\csname res@s1_isoright\endcsname{0.85}
\expandafter\def\csname res@s1_isowrong\endcsname{0.11}
\expandafter\def\csname res@s1_isowrong_k\endcsname{20}
\expandafter\def\csname res@s1_isowrong_n\endcsname{180}
\expandafter\def\csname res@s1_pickeach\endcsname{-0.05}
\expandafter\def\csname res@s1_pickeach_ci\endcsname{[-0.17, +0.07]}
\expandafter\def\csname res@s1_pickiso\endcsname{-0.03}
\expandafter\def\csname res@s1_pickiso_ci\endcsname{[-0.15, +0.08]}
\expandafter\def\csname res@s1_solve\endcsname{0.73}
\expandafter\def\csname res@s1_ta_each\endcsname{0.32}
\expandafter\def\csname res@s1_ta_iso\endcsname{0.30}
\expandafter\def\csname res@s1_ta_pick\endcsname{0.27}
\expandafter\def\csname res@s1_tp_iso_false\endcsname{0.30}
\expandafter\def\csname res@s1_tp_pick_false\endcsname{0.08}
\expandafter\def\csname res@s1_wrong_appear\endcsname{300}
\expandafter\def\csname res@s2_auroc_max\endcsname{0.98}
\expandafter\def\csname res@s2_auroc_min\endcsname{0.75}
\expandafter\def\csname res@s2_n_big\endcsname{10}
\expandafter\def\csname res@s2_unan_wrong_max\endcsname{0.06}
\expandafter\def\csname res@s3_auroc\endcsname{0.56}
\expandafter\def\csname res@s3_auroc_ci\endcsname{[+0.47, +0.65]}
\expandafter\def\csname res@s3_ta_unan\endcsname{0.75}
\expandafter\def\csname res@s3_ta_unique\endcsname{0.90}
\expandafter\def\csname res@s4_right_checker\endcsname{-0.12}
\expandafter\def\csname res@s4_right_checker_ci\endcsname{[-0.23, -0.02]}
\expandafter\def\csname res@s4_right_student\endcsname{+0.00}
\expandafter\def\csname res@s4_right_student_ci\endcsname{[-0.08, +0.08]}
\expandafter\def\csname res@s4_wrong_checker\endcsname{-0.08}
\expandafter\def\csname res@s4_wrong_checker_ci\endcsname{[-0.18, +0.02]}
\expandafter\def\csname res@s4_wrong_student\endcsname{-0.10}
\expandafter\def\csname res@s4_wrong_student_ci\endcsname{[-0.22, +0.02]}
\expandafter\def\csname res@s5_haiku_agent\endcsname{0.49}
\expandafter\def\csname res@s5_haiku_agentnorm\endcsname{0.58}
\expandafter\def\csname res@s5_haiku_handnorm\endcsname{0.43}
\expandafter\def\csname res@s5_haiku_step\endcsname{0.34}
\expandafter\def\csname res@s5_hand_majority\endcsname{0.53}
\expandafter\def\csname res@s5_sonnet_agent\endcsname{0.55}
\expandafter\def\csname res@s5_sonnet_agentnorm\endcsname{0.61}
\expandafter\def\csname res@s5_sonnet_handnorm\endcsname{0.57}
\expandafter\def\csname res@s5_sonnet_step\endcsname{0.40}
\expandafter\def\csname res@s6_ci\endcsname{[+0.03, +0.25]}
\expandafter\def\csname res@s6_diff\endcsname{+0.13}
\expandafter\def\csname res@s6_picked\endcsname{0.20}
\expandafter\def\csname res@s6_unpicked\endcsname{0.07}
\expandafter\def\csname res@s8_ta_neither\endcsname{0.68}
\expandafter\def\csname res@s8_tp_right\endcsname{0.68}
\expandafter\def\csname res@s9_adopt\endcsname{0.20}
\expandafter\def\csname res@s9_wrong_ci\endcsname{[-0.24, -0.02]}
\expandafter\def\csname res@s9_wrong_diff\endcsname{-0.12}
\expandafter\def\csname res@sb_conf_solved_ci\endcsname{[-0.19, +0.19]}
\expandafter\def\csname res@sb_conf_solved_d\endcsname{+0.00}
\expandafter\def\csname res@sb_conf_solved_n\endcsname{21}
\expandafter\def\csname res@sb_conf_unsolved_ci\endcsname{[+0.00, +0.00]}
\expandafter\def\csname res@sb_conf_unsolved_d\endcsname{+0.00}
\expandafter\def\csname res@sb_conf_unsolved_n\endcsname{3}
\expandafter\def\csname res@sb_repl_solved_ci\endcsname{[+0.00, +0.62]}
\expandafter\def\csname res@sb_repl_solved_d\endcsname{+0.25}
\expandafter\def\csname res@sb_repl_solved_n\endcsname{8}
\expandafter\def\csname res@sb_repl_unsolved_ci\endcsname{[+0.03, +0.26]}
\expandafter\def\csname res@sb_repl_unsolved_d\endcsname{+0.14}
\expandafter\def\csname res@sb_repl_unsolved_n\endcsname{70}
\expandafter\def\csname res@sb_screen_solved_ci\endcsname{[-0.23, +0.00]}
\expandafter\def\csname res@sb_screen_solved_d\endcsname{-0.11}
\expandafter\def\csname res@sb_screen_solved_n\endcsname{44}
\expandafter\def\csname res@sb_screen_unsolved_ci\endcsname{[-0.06, +0.44]}
\expandafter\def\csname res@sb_screen_unsolved_d\endcsname{+0.19}
\expandafter\def\csname res@sb_screen_unsolved_n\endcsname{16}

\newcommand{\res}[1]{\csname res@#1\endcsname}

\title{Beyond Its Competence, an LLM Judge Falls for Target-Absent Lineups: A Pre-registered Study of claude-haiku-5-5 on Officially Graded Math Solutions}
\author{Chaitanya Sharma\\Independent researcher}
\date{November 2026}

\begin{document}
\maketitle

\begin{abstract}
Selection steps in multi-agent pipelines often face pools in which no candidate is correct. Eyewitness research predicts that a judge shown such a target-absent lineup chooses the best-looking wrong candidate more often than it would accept any candidate judged alone. We test this for claude-haiku-5-5 on pools of three officially graded wrong solutions that models evaluated by MathArena produced for recent competition and arXiv-derived problems, asking the judge to choose the correct solution or NONE. In a pre-registered study on \res{c3_n_ta} fresh, independently locked problems, stratified before any judging by whether haiku solved the problem in two attempts, the pooled format added false choices on the \res{c3_n_unsolved} problems it could not solve: it chose a wrong solution in \res{c3h_unsolved_pick} of pools, while isolated verification accepted at least one of the same solutions in \res{c3h_unsolved_iso} (difference \res{c3h_du}, \res{lvl95} CI \res{c3h_du_ci}, Holm-confirmed). On those problems, a forced choice selected the wrong solution that haiku's own isolated verification accepts most (\res{c3h_h2u} above the pool's other wrong solutions, Holm-confirmed). The pre-registered primary test, that the penalty is larger on unsolved than on solved problems after adjusting for competition, was inconclusive (\res{c3h_b}, \res{c3h_b_ci}), because unsolved problems came almost entirely from one source. A pre-registered repeat with claude-sonnet-5-5 as the judge on the same strata showed the same two confirmed effects, and verifying haiku's forced picks with sonnet instead of haiku did not reduce the selection effect (\res{h3_haiku_d}, \res{h3_haiku_ci}). An earlier confirmation attempt failed for reasons we document. Seven exploratory screens, including a failure-attribution retest on Who\&When, are reported as screens. All results are scoped to the named model under stated settings. The ledger holds \res{ledger_total} hypotheses, including killed and null ones.
\end{abstract}

\section{Introduction}
Multi-agent systems for reasoning and discovery generate candidates and then select among them: AI co-scientist pipelines rank hypotheses in pairwise tournaments \citep{gottweis2025towards}, best-of-$n$ systems keep the candidate a verifier prefers \citep{cobbe2021training,brown2024large,stroebl2026limits}, and LLM judges score or compare outputs \citep{zheng2023judging}. Early in a hard search, the pool often contains no correct candidate. A selection step that cannot say ``none of these'' then passes a wrong answer downstream.

Eyewitness research has studied exactly this situation for decades. Witnesses shown a lineup that does not contain the culprit often choose the most similar person (relative judgment, \citealp{wells1984psychology}); instructions that the culprit may be absent reduce such choices \citep{malpass1981eyewitness}; and sequential presentation trades hits for fewer false identifications \citep{steblay2001eyewitness,steblay2011seventytwo}. LLM work has reported related symptoms: pairwise comparison can amplify judge biases \citep{jeong2025comparative}, and multiple-choice answering can behave like picking the least incorrect option \citep{wang2025llms}. Earlier unpublished pilots in this project, run on older models and synthetic items, found that a judge which rejected wrong candidates alone chose one of them from a target-absent pool far more often. Those pilots are lineage, not evidence for current models.

We ask whether the lineup effect holds for current judges on items that are hard for them and built from current models' real errors, and where selection without a rejection channel still fails.

\paragraph{Contributions.}
(1) A pre-registered, stratified test (C3) of the target-absent lineup effect for claude-haiku-5-5, with officially graded wrong solutions as candidates, solvability measured before judging by two independently written graders, and a single gate rule fixed in advance. On problems the judge cannot solve the lineup penalty is confirmed; whether it is specific to such problems is not established. (2) A confirmed failure of selection without a rejection channel on those problems: a forced pick selects exactly the wrong solution the same judge's verification passes most; and a screen showing that pairwise tournaments carry little signal that no candidate is right. (3) An honest record of a failed earlier confirmation (C2) and of what fixed it: per-model calibration, a gate-combination rule stated in advance, and two graders that must agree. (4) Seven exploratory screens on other multi-agent mechanisms, including a retest of failure attribution on Who\&When \citep{zhang2025which}, reported with every null. (5) A complete artifact: pre-registrations committed before data, raw outputs, and a ledger of every hypothesis.

\section{Related work}
\textbf{Selection under imperfect verification.} Resampling with an imperfect verifier has a false-positive ceiling \citep{stroebl2026limits}; coverage grows with samples \citep{brown2024large}; verifiers trained on math rank candidates \citep{cobbe2021training,lightman2024lets}. These analyses treat the verifier's error rates as fixed; we test whether presentation changes them.
\textbf{None of the above and abstention.} Language models can partly evaluate their own answers, including with none-of-the-above options \citep{kadavath2022language}, and abstention is surveyed by \citet{wen2024know} (TACL). \citet{wang2025llms} (COLING 2025) show that models often accept individual incorrect options when each is verified alone as true or false, relate that isolated verification to which option the model then chooses (on questions where a correct option is present), and find that models often still choose an option when the correct one is replaced and a no-answer option is offered. They do not test, within all-wrong pools, whether the forced choice is the candidate that isolated verification accepts most often, nor compare pooled choosing with a NONE option against isolated acceptance on problems the judge cannot solve; those are our contrasts, and their analysis linking isolated verification to choice is the closest prior result to our select-then-verify observation. \textbf{Comparative and listwise judging.} \citet{jeong2025comparative} (workshop) report that pairwise comparisons amplify biased preferences; \citet{wang2025llms} argue models answer multiple choice by selecting the least incorrect option; \citet{ji2026candidate} (preprint) study how candidate supply shapes the value of LLM judging but exclude all-wrong pools from their stratified and ranking analyses. We study target-absent pools directly, with and without a rejection option.
Concurrent work argues that what matters in verifier-guided best-of-$n$ is the reliability of the selected tail rather than average verifier accuracy \citep{zhang2026selected} (NeurIPS 2026 poster; we could read only its abstract). \textbf{Scalable oversight.} Debate can help judges \citep{irving2018ai,khan2024debating,kenton2024scalable}; we screen debate where both sides are wrong. \textbf{Failure attribution.} Who\&When provides human-labelled decisive agents and steps for failed multi-agent logs \citep{zhang2025which}. \textbf{Agreement as confidence.} Self-consistency uses majority agreement \citep{wang2023self}; we screen whether agreement stays informative on research-level problems.

\section{Method}
\subsection{Benchmarks and pools}
\textbf{MathArena} \citep{balunovic2025matharena} publishes recent competitions with official answers and the officially graded outputs of many models. A candidate is one output of a model in MathArena's 2026 lineup (which includes frontier models and smaller open models such as Qwen3-4B) whose final message has at least 300 characters; it is wrong or right by MathArena's grading, never by an LLM judge. A target-absent (TA) pool holds the first three distinct wrong answers for a problem; a target-present (TP) pool replaces one of them with a right candidate. Splits: calibration (AIME and HMMT February 2025), screens (SMT, CMIMC, BRUMO, APEX 2025, HMMT November 2025, arXiv-math June 2026), and held out for confirmation (HMMT February and AIME 2026; arXiv-math August and September 2026, whose candidates come from all models MathArena ran on those sets). Held-out pools were built and hash-locked by an independent agent before any confirmatory call. \textbf{Who\&When} \citep{zhang2025which}: 184 failed logs with human labels.

\subsection{Subjects}
Subjects are claude-haiku-5-5 and claude-sonnet-5-5 (resolved identifiers recorded on every call), called through Claude Code version \res{cli_version} in print mode with a fixed one-sentence system prompt, no tools, no project settings and an empty working directory; extended thinking is adaptive and on by default in this interface. Subjects saw only the task prompt; prompts contain no study context. The judge prompt also showed MathArena's grading of each candidate. Calls were made on 3 and 4 November 2026 (UTC). No claim extends to other versions, settings or labs.

\subsection{Arms}
ISO judges each candidate alone (ACCEPT or REJECT). PICK shows the pool and asks which solution is correct, if any, noting that none may be. EACH asks for a verdict on every candidate in one prompt. FORCED asks which is most likely correct with no NONE option. PAIR shows two candidates and forces a winner; a tournament runs every ordered pair. The TA choosing rate is the share of TA pools in which a wrong candidate is chosen (PICK) or at least one is accepted (ISO across its three calls, EACH). ISO uses three calls per pool and PICK one; we report this compute difference rather than equalize it.

\subsection{Calibration, inference and multiplicity}
A bank is interpreted for a model only if the model's solve rate on its problems is inside the pre-registered band (\res{band_lo} to \res{band_hi}) and its isolated false-accept rate is above zero. On the 2025 calibration split haiku solved \res{cal_haiku} and sonnet \res{cal_sonnet}, above the band, so that split is not interpreted; screens use the solve rate on their own problems. The problem is the unit of analysis; confidence intervals are percentile bootstraps over problems (TA and TP pools of one problem move together). Screens are exploratory. Two confirmatory studies were pre-registered: C2, which could not be interpreted (Section~\ref{sec:c2}), and C3 (Section~\ref{sec:c3}); each corrects its own hypothesis family with Holm's method. Every hypothesis tested or killed is in the ledger.

\subsection{The stratified confirmation (C3)}
C3 used only MathArena problems never shown to a subject in this project (\res{c3_n_ta} target-absent pools from 13 competitions, including six monthly arXiv-derived sets from December 2025 to May 2026), built and hash-locked by an independent agent; candidates are officially graded outputs of every model MathArena ran on that competition. Before any judging, haiku solved each problem twice; each reply was graded by our grader and by a second grader written by a separate agent that did not see the first, and the \res{c3_disagree} disagreements were resolved by hand against the official answer (all were equivalent forms) and committed before any judging call. Problems were SOLVED (both attempts correct, \res{c3_n_solved}), UNSOLVED (both wrong, \res{c3_n_unsolved}) or MIXED (\res{c3_n_mixed}, not judged). On SOLVED and UNSOLVED problems we ran ISO, PICK and FORCED. One gate, fixed in advance, governed interpretation: isolated-verdict parse rate at least \res{c3_gate_parse}, isolated acceptance of right candidates exceeding that of wrong candidates by at least \res{c3_gate_diff}, isolated false acceptance above zero, and at least 40 problems per stratum. The primary test (H1) is the coefficient of UNSOLVED in a linear model of the per-problem difference d (PICK chose a wrong candidate minus any isolated acceptance) with competition fixed effects and HC3 errors; secondary tests are mean d on UNSOLVED (H1u) and the select-then-verify difference within each stratum (H2s, H2u); the four form one Holm family. Spending was capped per arm before running.

\section{Results}
\subsection{A floor on stale items}
On 40 MATH problems with human-verified wrong GPT-4 solutions from PRM800K \citep{lightman2024lets,hendrycks2021measuring}, haiku, sonnet and opus never chose or accepted a wrong candidate in any format (isolated false accept \res{p1_haiku_isowrong}; solve rates \res{p1_haiku_solve}, \res{p1_sonnet_solve}, \res{p1_opus_solve}). These public items sit at the floor and test nothing about current judges; we report them because they were run. Outdated small open models were also run on these pools; under our rules they carry no weight and are listed only in the artifact.

\subsection{Screen: no lineup penalty on problems the judge mostly solves}
On 60 screen problems haiku was calibrated (solve rate \res{s1_solve}; isolated accept of wrong candidates \res{s1_isowrong}, that is \res{s1_isowrong_k} of the \res{s1_isowrong_n} distinct wrong candidates, each judged once although some appear in both a TA and a TP pool (\res{s1_wrong_appear} appearances); of right candidates \res{s1_isoright}). The TA choosing rate was \res{s1_ta_iso} for ISO, \res{s1_ta_pick} for PICK and \res{s1_ta_each} for EACH (Figure~\ref{fig:lineup}). PICK minus ISO was \res{s1_pickiso} \res{s1_pickiso_ci}; PICK minus EACH \res{s1_pickeach} \res{s1_pickeach_ci}. With a right candidate present, PICK chose a wrong one in \res{s1_tp_pick_false} of pools, against \res{s1_tp_iso_false} in which isolated verification accepted at least one wrong candidate; this comparison favours PICK by construction (PICK can choose at most one candidate, ISO-any can accept several) and is descriptive only.

\begin{figure*}[t]
\centering
\includegraphics[width=\textwidth]{figures/lineup}
\caption{Share of target-absent pools in which the judge chose (PICK) or accepted (ISO, any of three isolated calls; EACH) a wrong candidate, with Wilson \res{lvl95} intervals. Panels: screen and pre-registered held-out sets.}
\label{fig:lineup}
\end{figure*}

\subsection{An earlier confirmation that could not be interpreted (C2)}\label{sec:c2}
C2 pre-registered H1 (equivalence of PICK and ISO-any within $\pm$\res{eq_margin}), H2 (select then verify, below) and a calibration gate ``per model and set''. A review found that our answer parser dropped a leading LaTeX dollar and backslash, so some correct solutions were graded wrong; with the parser fixed (exps/c2\_confirm/AMENDMENTS.md A3), haiku solved \res{c2_haiku_conf_solve} of the 2026 contest problems, above the band, and \res{c2_haiku_repl_solve} of the arXiv-derived problems, below it. The union of both sets falls inside the band only by pooling (\res{c2_haiku_all_solve}); following our conservative ruling for such conflicts (made after the data, AMENDMENTS A2), the union is not interpretable either. C2 therefore tests nothing for haiku. For completeness: on the contest set H1 gave \res{c2_haiku_conf_h1d} \res{c2_haiku_conf_h1ci90} and H2 \res{c2_haiku_conf_h2}; on the union H1 \res{c2_haiku_all_h1d} \res{c2_haiku_all_h1ci90} and H2 \res{c2_haiku_all_h2} \res{c2_haiku_all_h2ci}; on the arXiv-derived set H1 \res{c2_haiku_repl_h1d} and H2 \res{c2_repl_h2}. Sonnet failed its gate (isolated-verdict parse rate \res{c2_sonnet_isoparse}, below the pre-registered threshold).

A post hoc split over the screen and C2 problems suggested the penalty sits on problems the judge cannot solve (\res{ph_unsolved_d} versus \res{ph_solved_d}), but solvability then rested on one attempt and one grader with a known bug, and was confounded with source. C3 was designed to test this properly.

\subsection{Stratified confirmation (C3)}\label{sec:c3}
Haiku passed the C3 gate (\res{c3h_gate}; isolated-verdict parse rate \res{c3h_isoparse}; isolated acceptance of wrong candidates \res{c3h_isowrong}, of right candidates \res{c3h_isoright}). Figure~\ref{fig:c3} shows the main result. On the \res{c3h_du_n} UNSOLVED problems, PICK chose a wrong solution in \res{c3h_unsolved_pick} of pools and isolated verification accepted at least one in \res{c3h_unsolved_iso}: d = \res{c3h_du} \res{c3h_du_ci}, Holm-adjusted $p$ \res{c3h_holm_H1u}, so H1u is \res{c3h_verdict_H1u}. On the \res{c3h_ds_n} SOLVED problems the rates were \res{c3h_solved_pick} and \res{c3h_solved_iso} (d = \res{c3h_ds} \res{c3h_ds_ci}, not a pre-registered test). The primary test H1, whether the penalty is larger on UNSOLVED problems after competition fixed effects, gave \res{c3h_b} \res{c3h_b_ci} (one-sided $p$ \res{c3h_b_p}, Holm \res{c3h_holm_H1}): \res{c3h_verdict_H1}. The unadjusted difference was \res{c3h_unadj} \res{c3h_unadj_ci}; the adjustment costs power because the contest competitions contributed only \res{c3h_fam_contest_nu} of \res{c3h_fam_contest_n} problems to the UNSOLVED stratum (within the arXiv-derived family alone, \res{c3h_fam_arxiv} \res{c3h_fam_arxiv_ci}). Select then verify: on UNSOLVED problems, the wrong solution FORCED picked was accepted by isolated verification \res{c3h_h2u} \res{c3h_h2u_ci} more often than the pool's other wrong solutions (Holm \res{c3h_holm_H2u}, \res{c3h_verdict_H2u}); on SOLVED problems the difference was \res{c3h_h2s} \res{c3h_h2s_ci} (\res{c3h_verdict_H2s}). Haiku's C3 judging cost \res{c3h_cost} US dollars.

\begin{figure*}[t]
\centering
\includegraphics[width=\textwidth]{figures/c3}
\caption{C3. Left of each pair: share of target-absent pools in which the judge chose (PICK) or accepted (ISO-any) a wrong solution, by haiku-solvability stratum, Wilson \res{lvl95} intervals. Right: isolated acceptance of the wrong solution a forced choice picked minus the mean for the other two, bootstrap \res{lvl95} intervals.}
\label{fig:c3}
\end{figure*}

\textbf{A second judge and cross-model verification.} A pre-registered amendment repeated C3 with claude-sonnet-5-5 as the judge on the same problems and strata (strata are defined by haiku, so for sonnet they measure the other model's competence), after a small pilot on items outside C3 fixed sonnet's parse failures in C2 by adding ``You have no tools and cannot run code'' to its system prompt (the C2 failures were plain-text imitations of tool calls). Sonnet passed the gate (parse rate \res{c3s_isoparse}). The pattern repeated: on haiku-UNSOLVED problems PICK chose a wrong solution in \res{c3s_unsolved_pick} of pools against \res{c3s_unsolved_iso} for ISO-any (d = \res{c3s_du} \res{c3s_du_ci}, Holm \res{c3s_holm_H1u}), the forced-pick difference was \res{c3s_h2u} \res{c3s_h2u_ci} (Holm \res{c3s_holm_H2u}), and the primary interaction was \res{c3s_b} \res{c3s_b_ci} (\res{c3s_verdict_H1}). H3 asked whether verifying haiku's forced pick with sonnet instead of haiku removes the select-then-verify inflation. It did not detectably: the inflation under same-model verification minus that under cross-model verification was \res{h3_haiku_d} \res{h3_haiku_ci} (one-sided $p$ \res{h3_haiku_p}; \res{h3_verdict}), and sonnet accepted haiku's forced pick in \res{h3_haiku_acc_sonnet} of pools against \res{h3_haiku_acc_haiku} for haiku itself. Within this pair of Claude models, the wrong solution a forced choice selects is also the one the other model is most likely to pass. Sonnet's C3 judging cost \res{c3s_cost} US dollars.

\subsection{Where selection fails: no rejection channel}
\textbf{Tournaments.} In a round-robin of forced pairwise comparisons (haiku, 40 problems), a unique winner emerged in \res{s3_ta_unique} of TA pools and a unanimous one in \res{s3_ta_unan}. The tournament's decisiveness separated TP from TA pools with AUROC \res{s3_auroc} \res{s3_auroc_ci}, so its own statistics give little warning that every candidate is wrong. A pilot on the stale items gave AUROC \res{p2_haiku_auroc} (haiku) and \res{p2_sonnet_auroc} (sonnet), even though those judges rejected every such candidate in isolation.
\textbf{Select, then verify.} In TA pools, the wrong candidate chosen by FORCED was accepted by haiku's isolated verification in \res{s6_picked} of pools, against \res{s6_unpicked} for the pool's other wrong candidates (difference \res{s6_diff}, \res{s6_ci}). This screen motivated H2, which C3 confirmed on problems the judge cannot solve (Section~\ref{sec:c3}). There, ranking first and then verifying with the same model sends the most dangerous candidate to the check most likely to pass it.

\subsection{Exploratory screens}
These are screens, never findings; each had a written kill rule. \textbf{Failure attribution (Who\&When).} Agent-level accuracy was \res{s5_haiku_agent} (haiku) and \res{s5_sonnet_agent} (sonnet) with strict name matching, \res{s5_haiku_agentnorm} and \res{s5_sonnet_agentnorm} after a post hoc name normalization; step-level accuracy \res{s5_haiku_step} and \res{s5_sonnet_step}. On hand-crafted logs, always naming the most frequent agent scores \res{s5_hand_majority}; normalized accuracy was \res{s5_haiku_handnorm} (haiku, below it) and \res{s5_sonnet_handnorm} (sonnet, close to it). \textbf{Critiques.} A false critique (``the final answer is wrong'') lowered acceptance of right candidates by \res{s4_right_checker} \res{s4_right_checker_ci} when attributed to an automated checker and by \res{s4_right_student} \res{s4_right_student_ci} when attributed to another student. \textbf{Debate with both sides wrong.} A haiku judge who saw only the debate answered NEITHER in \res{s8_ta_neither} of TA pairs and chose the right side in \res{s8_tp_right} of TP pairs. \textbf{Peer notes.} A wrong answer left by ``another agent'' changed haiku's accuracy by \res{s9_wrong_diff} \res{s9_wrong_ci}, adopted in \res{s9_adopt} of problems. \textbf{Agreement and diversity (MathArena's own runs, no subject calls).} For the \res{s2_n_big} models with at least 60 problems, the AUROC of self-consistency agreement ranged \res{s2_auroc_min} to \res{s2_auroc_max} and unanimous answers were wrong at most \res{s2_unan_wrong_max}; errors correlated less across different models (phi \res{s10_hetero}) than across samples of one model (\res{s10_homo}).

\begin{figure}[t]
\centering
\includegraphics[width=\columnwidth]{figures/whowhen}
\caption{Who\&When attribution accuracy (strict), agent (circles) and step (squares), Wilson \res{lvl95} intervals; dashed line: always naming the most frequent labelled agent (agent baseline); dotted line: a uniformly random step (step baseline).}
\label{fig:ww}
\end{figure}

\section{Threats to validity}
\textbf{Single vendor and version.} All subjects are Anthropic models; nothing here licenses statements about other labs or other Claude versions. \textbf{Interface settings.} Adaptive extended thinking was on by default and not controlled; with it on, an answer-only contamination probe was uninformative (answer-only accuracy \res{cal_haiku_nowork} and \res{cal_sonnet_nowork}), so contamination of 2025 to 2026 contests cannot be excluded. \textbf{Items.} Math final-answer tasks only; pools are restricted to problems with at least three distinct wrong answers from any model in the lineup, which is the intended regime but not a random sample; candidates from small models may be easier to reject and may lower the isolated false-accept rate. \textbf{Compute.} ISO uses three calls per pool, PICK one. \textbf{Analysis choices.} The Who\&When name normalization was added after seeing the data and is secondary. \textbf{Analysis errors found in review.} An answer-parsing bug, found by a reviewer after the first two review rounds, changed every solve rate; all numbers here use the fixed parser. \textbf{Stratification.} C3 strata are defined by two attempts of the judge, a noisy measure of competence; MIXED problems were not judged. UNSOLVED problems came almost entirely from arXiv-derived sets, so competence and source are entangled, which is why the source-adjusted primary test was inconclusive. \textbf{C2.} The C2 contest set has 24 problems; the pre-registration did not specify how per-set and union gates combine, and we resolved it conservatively after seeing the data. Sonnet's arm cost more than its pre-registered cap and failed its gate. One sonnet call was blocked by an API safeguard and one haiku solution call timed out; both are reported as missing. Two screens deviated from their written plans: the debate screen ran no consultancy baseline, and the peer-note screen's control was the judge's own unaided solution call from the lineup screen. \textbf{Reviewers.} All automated reviewers of this paper were Claude models and share blind spots with the subjects.

\section{Conclusion}
For claude-haiku-5-5, the eyewitness lineup effect appears where the judge is beyond its competence: on problems it could not solve, showing three wrong solutions together with a NONE option produced a wrong choice far more often than isolated verification produced a wrong acceptance, and a forced choice landed on exactly the wrong solution its own verification passes most. Whether the effect is specific to such problems, rather than present but small everywhere, is not established, because solvability and item source were entangled. For builders the reading is cautious and concrete: on problems the judge cannot itself solve, do not trust pooled selection to say ``none of these'', and do not assume that a second model from the same family will catch the pick: in our test, verification by claude-sonnet-5-5 removed none of the inflation. The next experiment is a stratified test in which solved and unsolved problems come from the same sources.

\section*{AI assistance and provenance}
The author set the research direction and constraints; an AI system (Claude, via Claude Code) wrote the rules and automated checks at the author's direction, ran the search, experiments, analysis and drafting under them, and automated Claude reviewers audited the draft. The author is responsible for the content. The ledger lists \res{ledger_total} hypotheses; subject calls cost about \res{cost_subjects} US dollars.

\bibliographystyle{plainnat}
\bibliography{refs}
\end{document}
